\subsection{Compatibilidad entre las publicaciones dodecafásica y analítica de alfa}
\label{subsec:alpha-limite-dos-vias}

La compatibilidad se establece a toda profundidad. La publicación
dodecafásica reversible y la representación analítica de la
coordenada de alfa
\(E_{21/22}\) son dos publicaciones internas con codominios distintos,
producidas a partir del mismo estado enriquecido APP--TRIT--TPK. El jet
\(E_{21/22}^{(9)}\) es un testigo finito exacto, pero no determina por sí solo
la cola. La completación que se construye después no se deduce del jet
desnudo: recibe la coordenada compartida del estado y publica su imagen
analítica mediante una recurrencia explícita. Por ello no constituye una
selección causal independiente; constituye una representación analítica
tipada del mismo punto generado.

Esta separación conserva el orden causal
\[
 \mathrm{APP}\longrightarrow\mathrm{TRIT}\longrightarrow\mathrm{TPK}
 \longrightarrow\mathcal X_{\alpha}^{\mathrm{enr}}
 \longrightarrow (P,E,\Phi,U_{12}^{\mathrm{sgn}},K)
 \longrightarrow\alpha_{\mathrm{HMT}}.
\]
Las coordenadas \(P,E,\Phi\) y el registro \(K\) son salidas anteriores del
mismo estado enriquecido. Su uso en la clausura de alfa es, por tanto, una
composición interna posterior y no una inyección de constantes
convencionales.

Para \(r\in6\mathbb N\), \(r\ge36\), sea
\(\widehat{\mathcal X}^{\mathrm{enr},\alpha}_r\) el sector de estados
truncados que prolongan la semilla interna de alfa fijada en \(R_{36}\). Se
escribe
\[
 x_r^\alpha=(\widetilde x_r,\chi_\alpha,
              \varepsilon_r^\alpha,N_r^\alpha),
\]
donde \(\chi_\alpha\) es el carácter interno del sector, no el valor escalar
publicado. Cada estado conserva residuo, cociente, acarreo, hoja, orientación,
fase, ruta, cilindro, frontera, firma, registro, memoria, supervivencia,
procedencia y profundidad. Si
\[
 d_r^\alpha=\operatorname{Dig}_{729}(\varepsilon_r^\alpha),\qquad
 \varepsilon_{r+6}^\alpha=729\varepsilon_r^\alpha-d_r^\alpha,
 \qquad N_{r+6}^\alpha=729N_r^\alpha+d_r^\alpha,
\]
la transición efectiva es
\[
 \mathcal U_r^\alpha=
 \operatorname{Upd}_r^\alpha\circ
 \operatorname{Tra}_r^\alpha\circ
 \operatorname{Sel}_r^\alpha,
 \qquad
 \operatorname{Ext}^{\alpha,\to}_{r+6,r}
 =\operatorname{graph}(\mathcal U_r^\alpha).
\]
La funcionalidad de estas extensiones, con la semilla \(R_{36}\) fijada,
produce una única sección compatible: la fase retorna mientras memoria y
profundidad avanzan. El dominio global de ambas determinaciones es
\begin{equation}
 \mathcal X_\alpha^{\mathrm{enr}}:=\widehat\Omega_\alpha
 :=\varprojlim_{r\ge36}
 \bigl(\widehat{\mathcal X}^{\mathrm{enr},\alpha}_r,
       \operatorname{Ext}^{\alpha,\to}_{r+6,r}\bigr).
 \label{eq:alpha-dominio-enriquecido-completo}
\end{equation}

\subsubsection{La clausura dodecafásica como sección completa}

Para \(N\geq1\), pongamos
\[
 \mathcal W_N:=\{0,\ldots,999\}^{N},
 \qquad
 \mathsf T_{\alpha,N}(\mathsf s):=(A_1,\ldots,A_N),
 \qquad \mathsf s\in\mathcal X_{\alpha}^{\mathrm{enr}},
\]
donde los bloques \(A_j\) se obtienen mediante
\eqref{eq:alpha-recurrencia} con la frontera compatible
\(c_{N+1}^{\infty}=\lfloor R_{N+1}(Z)\rfloor\). Si \(N'\geq N\), sea
\(\rho_{N',N}:\mathcal W_{N'}\to\mathcal W_N\) la restricción a los
primeros \(N\) bloques.

\begin{theorem}[Salida dodecafásica completa y compatibilidad proyectiva]
\label{thm:alpha-salida-dodecafasica-completa}
La familia \((\mathsf T_{\alpha,N})_{N\geq1}\) es compatible:
\begin{equation}
 \rho_{N',N}\circ\mathsf T_{\alpha,N'}
 =\mathsf T_{\alpha,N}
 \qquad(N'\geq N).
 \label{eq:alpha-via-a-proyectiva-exacta}
\end{equation}
Sus cilindros canónicos
\begin{equation}
 C_N^A(\mathsf s):=
 \left[
  \sum_{j=1}^{N}A_j1000^{-j},
  \sum_{j=1}^{N}A_j1000^{-j}+1000^{-N}
 \right)\cap I_\alpha
 \label{eq:alpha-cilindros-via-a-exactos}
\end{equation}
son anidados, tienen diámetro a lo sumo \(1000^{-N}\) y satisfacen
\begin{equation}
 \bigcap_{N\geq1}C_N^A(\mathsf s)
 =\{\alpha_A(\mathsf s)\},
 \qquad
 \alpha_A
 =\pi_{\mathrm{HMT}}+e_{\mathrm{HMT}}-\varphi_{\mathrm{HMT}}
  -4-\kappa_{\mathrm{per}}.
 \label{eq:alpha-salida-a-completa}
\end{equation}
En consecuencia, la primera vía publica una única sección infinita, que
denotamos \(\alpha_{\mathrm{HMT}}:=\alpha_A\).
\end{theorem}

\begin{proof}
La proposición~\ref{prop:alpha-normalizacion} proporciona la expansión
canónica única y la sucesión acotada única de acarreos. La identidad local
\eqref{eq:alpha-identidad-local} y su telescopado
\eqref{eq:alpha-telescopia} conservan la frontera remota; la proposición
\ref{ampl-alfa-prop-frontera} prueba que restringir una ejecución prolongada
recupera exactamente el prefijo anterior. Esto da
\eqref{eq:alpha-via-a-proyectiva-exacta}. La pertenencia de la suma completa
a cada cilindro se sigue de la identidad de colas, y
\(\operatorname{diam}C_N^A\leq1000^{-N}\to0\); la intersección contiene un
solo punto. Finalmente, \eqref{eq:alpha-identidad-completa} identifica ese
punto con la expresión de \eqref{eq:alpha-salida-a-completa}. Todos sus
coeficientes y signos se fijan antes de la publicación de alfa.
\end{proof}

La ventana \(\alpha_{12}^{(0)}\), cerrada con \(c_{13}=0\), sigue siendo un
cálculo finito exacto; no debe confundirse con la restricción de la sección
completa, cuya frontera compatible satisface \(c_{13}^{\infty}=-1\). En
particular,
\begin{equation}
 \rho_{36,12}(\alpha_{A,36})=\alpha_{A,12}
 \label{eq:alpha-rho-36-12-proyectiva}
\end{equation}
es una igualdad de palabras compatibles, no una igualdad entre palabras de
longitudes distintas ni una identificación de la ventana aislada con el
límite.

\subsubsection{Caracterización analítica exacta de orden nueve}

Sea
\[
 \mathbb K:=\mathbb Q\!\left(\pi_{\mathrm{HMT}},
                  \log_{10}\frac{10}{9}\right),
 \qquad
 E_9(x):=E_{21/22}^{(9)}(x)=P_9(x)-\delta .
\]
Las ecuaciones \eqref{eq:alpha-firma-coeficientes} y
\eqref{eq:alpha-t79} construyen \(P_9=P_6+T_{7:9}\) desde la firma del
estado y desde el resolvente nilpotente; la
proposición~\ref{prop:alpha-raiz-jet} demuestra que \(E_9\) posee una única
raíz simple \(\xi_9\in I_\alpha\), aislada por el intervalo racional
\eqref{eq:alpha-intervalo-jet}.

La bidireccionalidad finita puede escribirse sin introducir una cola
analítica. Definamos
\begin{align}
 B_9(x):={}&\delta+\frac74x^2-\frac{x^4}{20}
 +\frac{2x^5}{21}+\frac{x^6}{46}+\frac{x^7}{120}
 -\frac{x^8}{45}-\frac{2x^9}{495},
 \label{eq:alpha-b9-finito}\\
 q_x(T):={}&xT^2-B_9(x)T+x^3,\qquad x\in I_\alpha=(0,10^{-2}),
 \label{eq:alpha-q-involucion-finita}\\
 J_x:{\mathbb R}^{\times}\longrightarrow{\mathbb R}^{\times},
 \qquad J_x(T):={}&\frac{x^2}{T}.
\end{align}

\begin{proposition}[Equivalencia cuadrática e involución finitas]
\label{prop:alpha-involucion-finita-exacta}
Para todo \(x\in I_\alpha\) y todo \(T\in\mathbb R^\times\),
\begin{align}
 \frac{q_x(2\pi_{\mathrm{HMT}})}{2\pi_{\mathrm{HMT}}}
 &=P_9(x)-\delta=E_9(x),
 \label{eq:alpha-q-equivale-p9}\\
 J_x(J_x(T))&=T,
 \qquad
 T^2q_x(J_x(T))=x^2q_x(T).
 \label{eq:alpha-j-involucion-exacta}
\end{align}
Sea \(Z_x=\{T\in\mathbb R^\times:q_x(T)=0\}\). Por consiguiente,
\(J_x|_{Z_x}:Z_x\to Z_x\) es involutiva. Cuando \(E_9(x)=0\),
intercambia las raíces \(2\pi_{\mathrm{HMT}}\) y
\(x^2/(2\pi_{\mathrm{HMT}})\). Esta es una involución exacta sobre la fibra
escalar del jet de orden nueve; no actúa sobre el estado enriquecido.
\end{proposition}

\begin{proof}
Dividir \(q_x(T)\) por \(T\) y sustituir
\(T=2\pi_{\mathrm{HMT}}\) reproduce término a término la definición de
\(P_9(x)-\delta\). La primera identidad de
\eqref{eq:alpha-j-involucion-exacta} es inmediata. Para la segunda,
\[
 T^2q_x(x^2/T)
 =x^5-B_9(x)x^2T+x^3T^2=x^2q_x(T).
\]
Así, \(J_x\) preserva el conjunto de raíces y, por la fórmula de Viète,
intercambia las dos raíces cuyo producto es \(x^2\).
\end{proof}

El testigo decimal finito conserva entonces su fuerza exacta:
\begin{equation}
 \operatorname{Tr}_9(\xi_9^{-1})
 =\operatorname{Tr}_9\!\left((\alpha_{12}^{(0)})^{-1}\right)
 =137.035999084.
 \label{eq:alpha-testigo-finito-nueve}
\end{equation}
La igualdad \eqref{eq:alpha-testigo-finito-nueve} certifica la
caracterización finita y la involución precedente. No convierte la raíz de
un polinomio truncado en la raíz de toda prolongación compatible.

\subsubsection{El jet no determina por sí solo una prolongación analítica}

Sea
\[
 j_9:\mathbb K[[x]]\longrightarrow \mathbb K[x]/(x^{10})
\]
el truncamiento, y para \(h\in\mathbb K[x]\) pongamos
\begin{equation}
 E_h(x):=E_9(x)+x^{10}h(x).
 \label{eq:alpha-familia-prolongaciones-mismo-jet}
\end{equation}
Todos estos polinomios satisfacen \(j_9(E_h)=j_9(E_9)\).

\begin{theorem}[No unicidad de la raíz a partir del jet]
\label{thm:alpha-no-unicidad-prolongacion-jet}
No existe una aplicación
\[
 \mathfrak r_9:\mathbb K[x]/(x^{10})\longrightarrow I_\alpha
\]
que asigne, a partir del jet solamente, la raíz positiva de toda
prolongación polinómica de \(E_9\) que tenga una única raíz en
\(I_\alpha\). En particular, \(E_0\) y \(E_{1/729}\) poseen el mismo jet y
raíces positivas distintas.
\end{theorem}

\begin{proof}
La cota demostrada en la prueba de
\ref{prop:alpha-raiz-jet} da \(E_9'(x)>6.24\) en
\([0,10^{-2}]\). Para \(t\geq0\), escribamos
\(E_t(x)=E_9(x)+tx^{10}\). Entonces
\[
 E_t'(x)=E_9'(x)+10tx^9>6.24,
\]
mientras \(E_t(0)=-\delta<0\) y
\(E_t(10^{-2})>E_9(10^{-2})>0\). Cada \(E_t\) posee, por tanto, una raíz
positiva única \(\eta(t)\in I_\alpha\). Para \(t=1/729\),
\[
 E_{1/729}(\xi_9)=\frac{\xi_9^{10}}{729}>0;
\]
como \(E_{1/729}\) es estrictamente creciente y cambia de signo una sola
vez, su raíz satisface \(\eta(1/729)<\xi_9=\eta(0)\). Sin embargo,
\(j_9(E_{1/729})=j_9(E_0)\). Una aplicación que factorizase la raíz por
\(j_9\) tendría que asignar el mismo punto a ambos polinomios, contradicción.
\end{proof}

La variación no es sólo existencial. El teorema de la función implícita se
aplica porque \(\partial_xE_t(\eta(t))>6.24\), y proporciona
\begin{equation}
 \eta'(t)=-
 \frac{\eta(t)^{10}}
 {E_9'(\eta(t))+10t\eta(t)^9}<0.
 \label{eq:alpha-variacion-raiz-prolongacion}
\end{equation}
Así queda probado que los coeficientes posteriores al grado nueve modifican
la raíz aun cuando dejan invariante todo el jet construido hasta ese grado.
Esta conclusión no niega la cola producida por el estado enriquecido: prueba
únicamente que la raíz no factoriza por el truncamiento desnudo
\(j_9\). Escoger una cola \(h\) a partir de una raíz convencional externa
sería circular. La completación que sigue no consulta esa raíz ni una tabla
metrológica: evalúa directamente las coordenadas regionales y terminales del
estado. Como esa evaluación coincide algebraicamente con la publicación
entera, el resultado es una representación correlacionada, no una afirmación
de independencia entre dos selectores.

% Representación analítica correlacionada de alfa: acción, convergencia e
% intervalos.
\subsubsection{Representación analítica de la coordenada de alfa}
\label{subsec:alpha-lector-completo-rev08}

El control del jet anterior obliga a exhibir la cola: no basta nombrar un
lector completo. Sea \(\mathsf s\in\mathcal X_\alpha^{\mathrm{enr}}\). Antes
de la normalización decimal, el mismo estado ya publica la precoordenada
\begin{equation}
 a(\mathsf s):=p(\mathsf s)+e_0(\mathsf s)-f(\mathsf s)
                 -\kappa_{\mathrm{per}}(\mathsf s)\in I_\alpha,
 \label{eq:alpha-precoordenada-comun}
\end{equation}
con \(p=\pi_{\mathrm{HMT}}-3\),
\(e_0=e_{\mathrm{HMT}}-2\) y
\(f=\varphi_{\mathrm{HMT}}-1\). Equivalentemente,
\(a=\pi_{\mathrm{HMT}}+e_{\mathrm{HMT}}-
\varphi_{\mathrm{HMT}}-4-\kappa_{\mathrm{per}}\). No se introduce aquí
\(\alpha\) convencional: \(p,e_0,f\) son los tres caracteres regionales ya
generados y \(\kappa_{\mathrm{per}}\) es la lectura racional del registro
terminal. La clausura entera aplica a
\eqref{eq:alpha-precoordenada-comun} la normalización de
\eqref{eq:alpha-recurrencia}; la carta que sigue publica sobre la misma
precoordenada una representación analítica compatible. Su grafo de
dependencia recibe \((p,e_0,f,\kappa_{\mathrm{per}})\) directamente: no
contiene una arista \(\alpha_A\to\lambda\), aunque la igualdad demostrada más
abajo permita identificar después ambos nombres de publicación.

Escribamos \(q=1/729\), sea \(F_{\mathsf s}\) el cuerpo ordenado generado por
las coordenadas de \(\mathsf s\), y pongamos
\(E_{9,\mathsf s}(X)=E_{21/22}^{(9)}(\mathsf s;X)\). Definimos el coeficiente
de enlace
\begin{equation}
 \lambda(\mathsf s):=
 -\frac{E_{9,\mathsf s}(a(\mathsf s))
         \bigl(1-q\,a(\mathsf s)^3\bigr)}{a(\mathsf s)^{10}}.
 \label{eq:alpha-lambda-estado}
\end{equation}
Todos sus argumentos pertenecen al estado antes del reconocimiento
metrológico. En particular, \(\lambda\) no se obtiene ajustando una cadena
decimal objetivo. La ecuación~\eqref{eq:alpha-lambda-estado} debe leerse con
su tipo exacto: una vez que la primera publicación del estado ha producido
\(a(\mathsf s)\), fija la única prolongación de la clase geométrica que se
define a continuación. No constituye una segunda selección independiente de
\(a(\mathsf s)\), sino una representación analítica correlacionada del mismo
estado enriquecido.

En efecto, para \(\mu\in F_{\mathsf s}\) pongamos
\[
 E_{\mathsf s,\mu}(X)
 :=E_{9,\mathsf s}(X)+\mu\frac{X^{10}}{1-qX^3}.
\]
Como \(0<a(\mathsf s)<10^{-2}\) y \(qa(\mathsf s)^3<1\), se tiene
\begin{equation}
 E_{\mathsf s,\mu}(a(\mathsf s))=0
 \quad\Longleftrightarrow\quad
 \mu=\lambda(\mathsf s).
 \label{eq:alpha-unicidad-lambda-clase-geometrica}
\end{equation}
Así, la semilla de coeficientes no contiene libertad una vez fijados el estado
y la clase de prolongaciones \(q\)-geométricas. Esta unicidad es relativa a la
clase declarada; no afirma que el jet de orden nueve seleccione por sí solo una
cola entre todas las series analíticas posibles.

La acción local sobre bloques de tres coeficientes es
\begin{equation}
 \mathcal C_\alpha:F_{\mathsf s}^{3}\longrightarrow F_{\mathsf s}^{3},
 \qquad \mathcal C_\alpha(u,v,w)=q(u,v,w),
 \qquad c^{(0)}=(\lambda(\mathsf s),0,0).
 \label{eq:alpha-accion-coeficientes}
\end{equation}
Por tanto, para todo \(n\geq0\),
\begin{equation}
 b_{10+3n}=q^n\lambda(\mathsf s),\qquad
 b_{11+3n}=b_{12+3n}=0.
 \label{eq:alpha-recurrencia-coeficientes}
\end{equation}
Ésta es la regla que calcula cada coeficiente posterior al grado nueve. No
queda un parámetro libre en cada prolongación.

Para \(M\geq0\), sea
\begin{equation}
 E_{\mathsf s,M}(X):=E_{9,\mathsf s}(X)
 +\lambda(\mathsf s)\sum_{n=0}^{M}q^nX^{10+3n}.
 \label{eq:alpha-truncamientos-completos}
\end{equation}
La función completa es
\begin{equation}
 E_{\mathsf s,\infty}(X)
 :=E_{9,\mathsf s}(X)
 +\lambda(\mathsf s)\frac{X^{10}}{1-X^3/729}.
 \label{eq:alpha-funcion-completa-explicita}
\end{equation}

Para formular el sistema proyectivo sin tipos implícitos, fijamos
\(d_M=12+3M\) y definimos la fibración de jets
\begin{equation}
 \mathfrak J_M:=
 \coprod_{\mathsf s\in\mathcal X_\alpha^{\mathrm{enr}}}
 \{\mathsf s\}\times F_{\mathsf s}[X]/(X^{d_M+1}).
 \label{eq:alpha-espacio-jets-rev08}
\end{equation}
Si \(M'\ge M\), la restricción es
\begin{equation}
 \tau_{M',M}(\mathsf s,[P]_{d_{M'}})
 :=(\mathsf s,[P]_{d_M}).
 \label{eq:alpha-truncamiento-jets-rev08}
\end{equation}

\begin{proposition}[Naturalidad de la recurrencia analítica]
\label{prop:alpha-naturalidad-recurrencia-explicita}
Sea
\[
 \Gamma_M(\mathsf s)
 =\bigl(\mathsf s,[E_{\mathsf s,M}]_{d_M}\bigr).
\]
Si \(M'\geq M\), el truncamiento \(\tau_{M',M}\) satisface
\begin{equation}
 \tau_{M',M}\circ\Gamma_{M'}=\Gamma_M,
 \qquad
 \Gamma_{E21}(\mathsf s)=\varprojlim_M\Gamma_M(\mathsf s).
 \label{eq:alpha-naturalidad-explicita}
\end{equation}
La proyección sobre los tres grados nuevos es precisamente
\(\mathcal C_\alpha^{M+1}(c^{(0)})\).
\end{proposition}

\begin{proof}
La ecuación~\eqref{eq:alpha-recurrencia-coeficientes} fija el bloque de
grados \(10+3n,11+3n,12+3n\). Al pasar de \(M\) a \(M+1\) sólo se añade
el bloque siguiente; truncarlo recupera literalmente el polinomio anterior.
La compatibilidad para \(M'\geq M\) se obtiene por composición. Así, el
límite inverso no es una sección escogida del jet desnudo, sino la órbita
única de \(c^{(0)}\) bajo \(\mathcal C_\alpha\).
\end{proof}

\subsubsection{Convergencia uniforme, raíz simple e intervalos racionales}
\label{subsec:alpha-cotas-completas-rev08}

Los cilindros regionales y el registro periódico proporcionan cotas mediante
aritmética de intervalos racionales. Para cada coordenada
\[
 c\in\{\pi_{\mathrm{HMT}},e_{\mathrm{HMT}},\varphi_{\mathrm{HMT}}\},
\]
sea \(P_c\) su prefijo nonádico certificado de mil cifras. El dato utilizado
no es el racional \(P_c\) como si fuese el valor completo, sino el cilindro
semiabierto y su cierre exterior
\[
 \mathcal C_c^{(1000)}=[P_c,P_c+\varepsilon),
 \qquad
 \overline{\mathcal C}_c^{(1000)}=[P_c,P_c+\varepsilon],
 \qquad \varepsilon=10^{-1000}.
\]
La aritmética dirigida opera conservadoramente con esos cierres. Si
\(P_\pi,P_e,P_\varphi\) denotan los tres prefijos, se obtiene
\begin{equation}
\begin{aligned}
 a_-&:=P_\pi+P_e-(P_\varphi+\varepsilon)-4-
       \kappa_{\mathrm{per}},\\
 a_+&:=(P_\pi+\varepsilon)+(P_e+\varepsilon)-P_\varphi-4-
       \kappa_{\mathrm{per}},\\
 A_{1000}&:=[a_-,a_+],
 \qquad a(\mathsf s)\in[a_-,a_+]=A_{1000},
 \qquad \operatorname{diam}A_{1000}=3\varepsilon.
\end{aligned}
\label{eq:alpha-propagacion-colas-rev08}
\end{equation}
La extensión intervalar natural de
\eqref{eq:alpha-lambda-estado} transporta después \(A_{1000}\), el cilindro
de \(\pi_{\mathrm{HMT}}\) y la cota racional dirigida de \(\delta\) a un
intervalo cerrado \(\Lambda_{1000}\ni\lambda(\mathsf s)\); no se congela
ninguna de las tres colas. En particular,
\begin{equation}
\begin{aligned}
 \frac{7297352569283800997285105472380662}{10^{36}}
 &<a(\mathsf s)<
 \frac{7297352569283800997285105472380663}{10^{36}},\\
 -8.711\mathbin{\times}10^{-6}
 &<\lambda(\mathsf s)<-8.709\mathbin{\times}10^{-6}.
\end{aligned}
 \label{eq:alpha-cotas-precoordenada-lambda}
\end{equation}
Estas desigualdades usan prefijos certificados con su cota racional de
cola; no reciben una tabla de alfa. El verificador del paquete las reproduce
desde las tres salidas nonádicas y desde la carta terminal.

Sea
\[
 r:=\frac{10^{-6}}{729}<1.372\mathbin{\times}10^{-9}.
\]
Para \(0\leq x\leq10^{-2}\), la cola posterior a \(M\) satisface
\begin{equation}
 \left|E_{\mathsf s,\infty}(x)-E_{\mathsf s,M}(x)\right|
 \leq
 8.711\mathbin{\times}10^{-26}\,\frac{r^{M+1}}{1-r}.
 \label{eq:alpha-cota-cola-uniforme}
\end{equation}
Además,
\begin{equation}
 \left|\frac{\mathrm d}{\mathrm dx}
  \left(\lambda\frac{x^{10}}{1-x^3/729}\right)\right|
 \leq 8.711\mathbin{\times}10^{-6}
 \left(\frac{10\,10^{-18}}{1-r}
 +\frac{3\,10^{-24}/729}{(1-r)^2}\right)
 <8.712\mathbin{\times}10^{-23}.
 \label{eq:alpha-cota-derivada-cola}
\end{equation}
La cola de las derivadas posterior al bloque \(M\) admite, además, la
estimación dependiente de la profundidad
\begin{equation}
 \sup_{0\le x\le10^{-2}}
 \left|E'_{\mathsf s,\infty}(x)-E'_{\mathsf s,M}(x)\right|
 \le 8.711\mathbin{\times}10^{-24}\,r^{M+1}
 \left(\frac{13+3M}{1-r}+\frac{3r}{(1-r)^2}\right)
 \xrightarrow[M\to\infty]{}0.
 \label{eq:alpha-cota-resto-derivadas-rev08}
\end{equation}

\begin{theorem}[Función límite y raíz simple única de la carta correlacionada]
\label{thm:alpha-raiz-completa-explicita}
La sucesión \(E_{\mathsf s,M}\) converge uniformemente, junto con sus
derivadas, a \(E_{\mathsf s,\infty}\) en
\(\overline I_\alpha=[0,10^{-2}]\). Se cumple
\begin{equation}
 E_{\mathsf s,\infty}(a(\mathsf s))=0,
 \qquad
 \inf_{x\in\overline I_\alpha}
 E'_{\mathsf s,\infty}(x)>6.239.
 \label{eq:alpha-raiz-y-derivada-completas}
\end{equation}
Por consiguiente, \(a(\mathsf s)\) es la raíz simple única de la función
completa en \(I_\alpha\).
\end{theorem}

\begin{proof}
La razón de dos términos consecutivos de la cola está acotada por \(r<1\),
lo que da \eqref{eq:alpha-cota-cola-uniforme}; la derivación término a
término y la suma de las dos series geométricas dan
\eqref{eq:alpha-cota-derivada-cola} y
\eqref{eq:alpha-cota-resto-derivadas-rev08}. La prueba del jet establece
\(E'_{9,\mathsf s}>6.24\) en \(\overline I_\alpha\); por tanto, la
perturbación completa deja la derivada por encima de \(6.239\).
Finalmente, al sustituir \eqref{eq:alpha-lambda-estado} en
\eqref{eq:alpha-funcion-completa-explicita}, los dos sumandos evaluados en
\(a(\mathsf s)\) se cancelan exactamente. La función es estrictamente
creciente; su única raíz es, pues, simple y coincide con esa precoordenada.
\end{proof}

Para hacer efectiva la identificación a toda profundidad, las evaluaciones
dirigidas de los extremos certifican primero
\[
 E_{\mathsf s,\infty}(0)<0<
 E_{\mathsf s,\infty}(10^{-2}),
 \qquad
 E'_{\mathsf s,\infty}\geq\mu:=6239/1000>c:=6
 \quad\hbox{en }[0,10^{-2}].
\]
Por continuidad existe una raíz en el intervalo y, por la segunda desigualdad,
es única. Sólo después de fijar esta existencia, la monotonía y la invariante
de que la horquilla contiene la raíz se aplica el rescate por el teorema del
valor medio. Partimos de \(D_0=[0,10^{-2}]\). Si
\(D_j=[r_j,s_j]\), sea \(m_j=(r_j+s_j)/2\), y sea
\([F_j^-,F_j^+]\) una inclusión racional dirigida de
\(E_{\mathsf s,\infty}(m_j)\), obtenida propagando los cilindros de
\(\pi_{\mathrm{HMT}},e_{\mathrm{HMT}},\varphi_{\mathrm{HMT}}\),
\(\delta\) y \(\lambda\), con
\begin{equation}
 0\leq F_j^+-F_j^-\leq \frac{c}{4}\,(s_j-r_j).
 \label{eq:alpha-anchura-evaluacion-rescate-rev08}
\end{equation}
Cuando \(F_j^+<0\) se conserva la mitad derecha;
cuando \(F_j^->0\), la mitad izquierda. Si
\(F_j^-\leq0\leq F_j^+\), no se intenta decidir si el punto medio es
exactamente la raíz. Sea \(w_j=s_j-r_j\). La cota uniforme
\(E'_{\mathsf s,\infty}\geq c=6\) y el teorema del valor medio implican que
cualquier raíz contenida en \(D_j\) pertenece a
\begin{equation}
 D_{j+1}=D_j\cap
 \left[m_j-\frac{w_j}{4},
             m_j+\frac{w_j}{4}\right].
 \label{eq:alpha-rescate-derivada-rev08}
\end{equation}
En efecto, sea \(\xi_j\) la raíz y sea
\(y_j=E_{\mathsf s,\infty}(m_j)\in[F_j^-,F_j^+]\). Como el recinto contiene
cero y satisface \eqref{eq:alpha-anchura-evaluacion-rescate-rev08}, se tiene
\(|y_j|\leq cw_j/4\). El teorema del valor medio proporciona
\(|m_j-\xi_j|\leq |y_j|/c\leq w_j/4\), lo que prueba
\eqref{eq:alpha-rescate-derivada-rev08}. El argumento incluye el caso
\(y_j=0\): la raíz está entonces en el punto medio, pero el algoritmo no
necesita reconocer esa igualdad para conservarla.

Si el recinto disponible no satisface
\eqref{eq:alpha-anchura-evaluacion-rescate-rev08}, se refinan primero los
cilindros fuente y se repite la evaluación; no se declara un signo ni se
acepta una contracción insuficiente. En cualquiera de las tres ramas el nuevo
intervalo conserva la raíz y tiene diámetro a lo sumo \(w_j/2\). Para cada
profundidad objetivo finita, el refinamiento dirigido permite satisfacer la
cota de anchura sin comparar de forma no certificada un número real con cero.
La monotonía demuestra inductivamente
\begin{equation}
 a(\mathsf s)\in D_{j+1}\subseteq D_j,\qquad
 \operatorname{diam}D_{j+1}\leq
 \tfrac12\operatorname{diam}D_j,
 \qquad \operatorname{diam}D_j\longrightarrow0,
 \label{eq:alpha-biseccion-racional}
\end{equation}
donde la desigualdad es una cota de contracción y no una igualdad impuesta.
El certificado ejecutable aplica
esta regla a veinticuatro niveles en base mil y comprueba que el cierre
exterior completo \(A_{1000}\), no sólo su extremo inferior, queda dentro de
cada cilindro publicado. Sus pruebas negativas rechazan un recinto demasiado
ancho, un recinto desordenado, un punto que no sea el medio, un dominio sin
cambio de signo y una horquilla degenerada. La prolongación coinductiva permite
reemplazar mil por cualquier profundidad requerida.

Concretamente, para \(S_N=1000^N\) el verificador calcula
\begin{equation}
 j_N^-:=\lfloor S_Na_-\rfloor,
 \qquad j_N^+:=\lfloor S_Na_+\rfloor
 \label{eq:alpha-indices-cilindro-dirigidos-rev08}
\end{equation}
y sólo publica el nivel si \(j_N^-=j_N^+\). Esta igualdad, que incluye el
extremo superior del cierre exterior, prueba
\(A_{1000}\subset[j_N^-/S_N,(j_N^-+1)/S_N)\) y evita atribuir a la celda
izquierda un extremo situado exactamente sobre su frontera derecha.

Sea \(C_N^A(\mathsf s)\) el cilindro entero de
\eqref{eq:alpha-cilindros-via-a-exactos}. Elegimos recursivamente
\(m(N)>m(N-1)\) de modo que
\(\operatorname{diam}D_{m(N)}<1000^{-N}/4\), y definimos
\begin{equation}
I_{B,N}:=D_{m(N)}\cap C_N^A(\mathsf s).
 \label{eq:alpha-intervalos-b-completos}
\end{equation}
Definimos
\(\ell_N:=\inf I_{B,N}\) y \(u_N:=\sup I_{B,N}\). Ambos son racionales;
el intervalo es no vacío porque los dos factores contienen
\(a(\mathsf s)\). Puesto que el cilindro es semiabierto, \(u_N\) puede no
pertenecer a \(I_{B,N}\). La evaluación en ese extremo se entiende mediante
la extensión continua de \(E_{\mathsf s,\infty}\) al cierre exterior. Así,
\begin{equation}
 I_{B,N+1}\subseteq I_{B,N}\subseteq C_N^A(\mathsf s),
 \qquad \operatorname{diam}I_{B,N}\longrightarrow0,
 \qquad E_{\mathsf s,\infty}(\ell_N)\leq0\leq
 E_{\mathsf s,\infty}(u_N).
 \label{eq:alpha-inclusion-signos-completa}
\end{equation}
La cota~\eqref{eq:alpha-cota-cola-uniforme} permite escoger un
truncamiento que conserve el signo en cada extremo donde la función completa
no se anula. Si un extremo coincide con la raíz, se utiliza la identidad de
la función completa y la pertenencia ya demostrada; no se atribuye ese mismo
signo a sus truncamientos. En efecto, el resto exacto es
\[
 E_{\mathsf s,\infty}(x)-E_{\mathsf s,M}(x)
 =\lambda(\mathsf s)x^{10}
   \frac{(qx^3)^{M+1}}{1-qx^3}.
\]
En el dominio positivo considerado, donde \(\lambda(\mathsf s)<0\), se
obtiene en la raíz
\[
 E_{\mathsf s,M}(a(\mathsf s))
 =-\lambda(\mathsf s)a(\mathsf s)^{10}
   \frac{(qa(\mathsf s)^3)^{M+1}}{1-qa(\mathsf s)^3}>0
 \qquad(M<\infty).
\]
Por ello, la certificación de los signos débiles en
\eqref{eq:alpha-inclusion-signos-completa} se refiere a
\(E_{\mathsf s,\infty}\). Su evaluación utiliza la función completa o el
truncamiento acompañado del resto con signo. Esto cubre tanto el ínfimo
perteneciente como el supremo eventualmente excluido del cilindro.

\begin{proposition}[Cuadrados de compatibilidad y cilindros comunes]
\label{prop:alpha-cuadrados-compatibilidad}
Las restricciones de palabras y de publicaciones analíticas conmutan:
\[
\begin{array}{ccc}
 \mathcal X_\alpha^{\mathrm{enr}}&
 \xrightarrow{\ \mathsf T_{\alpha,N'}\ }&\mathcal W_{N'}\\[1mm]
 \Big\Vert&&\Big\downarrow\rho_{N',N}\\[1mm]
 \mathcal X_\alpha^{\mathrm{enr}}&
 \xrightarrow{\ \mathsf T_{\alpha,N}\ }&\mathcal W_N
\end{array}
\qquad
\begin{array}{ccc}
 \mathcal X_\alpha^{\mathrm{enr}}&
 \xrightarrow{\ \Gamma_{M'}\ }&\mathfrak J_{M'}\\[1mm]
 \Big\Vert&&\Big\downarrow\tau_{M',M}\\[1mm]
 \mathcal X_\alpha^{\mathrm{enr}}&
 \xrightarrow{\ \Gamma_M\ }&\mathfrak J_M .
\end{array}
\]
La familia \(I_{B,N}\) proporciona la conexión efectiva entre ambos
cuadrados y satisface materialmente
\(I_{B,N}\subseteq C_N^A\) para todo \(N\).
\end{proposition}

\begin{proof}
El primer cuadrado es la compatibilidad de la normalización entera; el segundo
es la proposición~\ref{prop:alpha-naturalidad-recurrencia-explicita}. La
inclusión se sigue de la definición~\eqref{eq:alpha-intervalos-b-completos},
y la no vacuidad y los signos se siguen del
teorema~\ref{thm:alpha-raiz-completa-explicita}. Ningún paso identifica la
raíz de un truncamiento con la de una cola arbitraria.
\end{proof}

Definimos ahora, sin dejar implícito el lector de prefijos,
\begin{equation}
 \alpha_B(\mathsf s):=
 \operatorname*{arg\,zero}_{x\in I_\alpha}E_{\mathsf s,\infty}(x),
 \qquad
 \operatorname{pref}_N(x):=
 \frac{\lfloor1000^Nx\rfloor}{1000^N},
 \qquad
 \alpha_{B,N}:=\operatorname{pref}_N(\alpha_B).
 \label{eq:alpha-definicion-b-y-prefijos-rev08}
\end{equation}

\begin{theorem}[Igualdad proyectiva de las dos publicaciones correlacionadas]
\label{prop:alpha-dos-limites}
Para todo \(N\), la vía entera y la vía analítica publican el mismo prefijo:
\begingroup
\renewcommand{\theequation}{A7\(_0\)}
\begin{equation}
 \boxed{\alpha_{A,N}=\operatorname{pref}_N(\alpha_A)
 =\operatorname{pref}_N(\alpha_B)=\alpha_{B,N}.}
 \label{eq:alpha-dos-vias-cada-nivel}
\end{equation}
\endgroup
En el límite,
\begingroup
\renewcommand{\theequation}{A8\(_0\)}
\begin{equation}
 \boxed{\alpha_A
 =\operatorname*{arg\,zero}_{x\in I_\alpha}
       E_{\mathsf s,\infty}(x)
 =\alpha_B=\alpha_{\mathrm{HMT}}.}
 \label{eq:alpha-dos-vias-limite-coinductivo}
\end{equation}
\endgroup
\end{theorem}

\begin{proof}
La normalización entera publica exactamente \(a(\mathsf s)\). El
teorema~\ref{thm:alpha-raiz-completa-explicita} demuestra que la carta
analítica representa ese mismo punto como raíz simple única. Los
cilindros semiabiertos \(C_N^A\) fijan un prefijo canónico y
\(I_{B,N}\subseteq C_N^A\) contiene \(\alpha_B=\alpha_A\). Por tanto, sus
prefijos coinciden para cada \(N\); el anidamiento y los diámetros tendentes a
cero dan la igualdad de los límites. El reconocimiento
metrológico ocurre después.
\end{proof}

La función analítica constituye una publicación correlacionada del mismo
estado, no una segunda derivación causalmente independiente de la clausura
entera. Esta precisión no reduce la igualdad demostrada: identifica su
mecanismo exacto, impide atribuir al jet una cola que no determina y evita la
promoción de un valor convencional a generador.

\endinput

% El bloque que sigue se conserva únicamente como antecedente editorial de
% REV04. \endinput impide que forme parte de la fuente activa o del PDF REV08.
\subsubsection{Lector completo, naturalidad y cilindros comunes}

Sea \(F_{\mathsf s}\) el cuerpo de coeficientes de la firma de
\(\mathsf s\in\mathcal X_\alpha^{\mathrm{enr}}\) y sea \(d_m=6+3m\). Se
definen
\[
 \mathfrak J_m:=\coprod_{\mathsf s\in\mathcal X_\alpha^{\mathrm{enr}}}
 \{\mathsf s\}\times F_{\mathsf s}[X]/(X^{d_m+1}),\qquad
 \Gamma_m(\mathsf s)=
 \bigl(\mathsf s,E_{21/22}^{(d_m)}(\mathsf s;X)\bigr).
\]
El lector completo registrado \(\Gamma_{E21}\) produce estas publicaciones.
Para \(m'\ge m\), el truncamiento
\(
 \tau_{m',m}(\mathsf s,[f]_{d_{m'}})
 =(\mathsf s,[f]_{d_m})
\)
satisface
\begin{equation}
 \tau_{m',m}\circ\Gamma_{m'}=\Gamma_m,
 \qquad
 E_{21/22}:=\varprojlim_mE_{21/22}^{(d_m)}.
 \tag{A3}
 \label{eq:alpha-sistema-analitico-completo}
\end{equation}
No se añade una sección del truncamiento desnudo. Si se explicita la
publicación de los tres coeficientes siguientes, ésta es la proyección
derivada
\[
 \Lambda_{d_m+1:d_m+3}
 :=\operatorname{pr}_{d_m+1:d_m+3}\circ
   \Gamma_{E21}\circ\mathcal U_{\mathrm{TPK}},
\]
no un generador independiente ni una lectura del valor objetivo.

Para un conjunto cofinal \(\mathcal N_\alpha\), el lector publica una función
cofinal \(m:\mathcal N_\alpha\to\mathbb N\) e intervalos
\(I_{B,N}=[\ell_N,u_N]\Subset I_\alpha\) tales que
\begin{equation}
\begin{gathered}
 I_{B,N'}\subseteq I_{B,N}\ (N'\ge N),\qquad
 \operatorname{diam}I_{B,N}\longrightarrow0,\qquad
 I_{B,N}\subseteq C_N^A,\\
 E_{21/22}^{(d_{m(N)})}(\ell_N)\le0\le
 E_{21/22}^{(d_{m(N)})}(u_N),\qquad
 \inf_{x\in I_{B,N}}\partial_xE_{21/22}^{(d_{m(N)})}(x)>0,\\
 E_{21/22}^{(d_{m(N)})}\longrightarrow E_{21/22}
 \quad\text{uniformemente en }I_\alpha .
\end{gathered}
 \tag{A4}
 \label{eq:alpha-via-b-proyectiva}
\end{equation}

\begin{proposition}[Cuadrados de compatibilidad y cilindros comunes]
\label{prop:alpha-cuadrados-compatibilidad}
Las restricciones de palabras y jets conmutan separadamente:
\[
\begin{array}{ccc}
 \mathcal X_\alpha^{\mathrm{enr}}&
 \xrightarrow{\ \mathsf T_{\alpha,N'}\ }&\mathcal W_{N'}\\[1mm]
 \Big\Vert&&\Big\downarrow\rho_{N',N}\\[1mm]
 \mathcal X_\alpha^{\mathrm{enr}}&
 \xrightarrow{\ \mathsf T_{\alpha,N}\ }&\mathcal W_N
\end{array}
\qquad
\begin{array}{ccc}
 \mathcal X_\alpha^{\mathrm{enr}}&
 \xrightarrow{\ \Gamma_{m'}\ }&\mathfrak J_{m'}\\[1mm]
 \Big\Vert&&\Big\downarrow\tau_{m',m}\\[1mm]
 \mathcal X_\alpha^{\mathrm{enr}}&
 \xrightarrow{\ \Gamma_m\ }&\mathfrak J_m .
\end{array}
 \tag{A5}
\]
La conexión entre ambos cuadrados es la inclusión certificada
\(I_{B,N}\subseteq C_N^A\), con anidamiento, contracción y cofinalidad. Por
ello, el punto único
\(x_B:=\bigcap_{N\in\mathcal N_\alpha}I_{B,N}\) satisface
\begin{equation}
 \forall N\in\mathcal N_\alpha,\qquad
 \operatorname{pref}_N(x_B)=\alpha_{A,N}=:\alpha_{B,N}.
 \tag{A6}
 \label{eq:alpha-igualdad-cada-nivel}
\end{equation}
\end{proposition}

\begin{proof}
La compatibilidad de la vía A es
\eqref{eq:alpha-via-a-proyectiva-exacta}; la de la vía B es
\eqref{eq:alpha-sistema-analitico-completo}. Cada jet posee una raíz simple
en su intervalo por el cambio de signo y la cota derivativa de
\eqref{eq:alpha-via-b-proyectiva}. La familia anidada determina \(x_B\), y la
convergencia uniforme da \(E_{21/22}(x_B)=0\). La inclusión cofinal en los
cilindros de la vía A fuerza sus prefijos. No se ha supuesto que truncar un
polinomio preserve su raíz.
\end{proof}

\begin{theorem}[Igualdad proyectiva y coinductiva de las dos vías]
\label{prop:alpha-dos-limites-historico-rev04}
Para todo nivel compatible del sistema cofinal común,
\begingroup
\renewcommand{\theequation}{A7}
\begin{equation}
 \boxed{\alpha_{A,N}=\rho_N(\alpha_A)
 =\rho_N(\alpha_B)=\alpha_{B,N}.}
 \label{eq:alpha-dos-vias-cada-nivel-historico-rev04}
\end{equation}
\endgroup
En el límite coinductivo,
\begingroup
\renewcommand{\theequation}{A8}
\begin{equation}
 \boxed{\alpha_A
 =\varprojlim_N\alpha_{A,N}
 =\varprojlim_N\alpha_{B,N}
 =\operatorname*{arg\,zero}_{x\in I_\alpha}E_{21/22}(x)
 =\alpha_B=\alpha_{\mathrm{HMT}}.}
 \label{eq:alpha-dos-vias-limite-coinductivo-historico-rev04}
\end{equation}
\endgroup
Ningún valor convencional interviene: el reconocimiento metrológico es
posterior a las dos publicaciones internas.
\end{theorem}

\begin{proof}
La vía A forma cilindros compatibles de diámetro \(1000^{-N}\); la vía B,
intervalos anidados de diámetro tendente a cero. Por
\eqref{eq:alpha-via-b-proyectiva}, el punto de la segunda familia es la raíz
simple y única del núcleo completo. Como pertenece a una subfamilia cofinal
de los cilindros de la primera, ambas intersecciones contienen el mismo punto.
La ecuación~\eqref{eq:alpha-igualdad-cada-nivel} da la igualdad a toda
profundidad compatible y el límite inverso da (A8). Las ejecuciones finitas
sólo certifican proyecciones; no acotan el enunciado.
\end{proof}

La raíz simple \(\xi_9\) y el testigo
\eqref{eq:alpha-testigo-finito-nueve} son verificaciones finitas de este
sistema. No definen el límite ni reducen (A7)--(A8) a una coincidencia
truncada: la igualdad completa procede de las extensiones del estado, de los
dos cuadrados de compatibilidad y de la unicidad de la intersección cofinal.
